% =============================================================================
%  Forecast targets and loss functions:
%  deep time-series models vs. the multi-horizon HAR-RV arrangement
%
%  Compiles standalone:      pdflatex dl_vs_har_targets_losses.tex
%  Or drop into a thesis:    copy everything between the two
%                            "%%% BODY BEGIN/END" markers into a chapter and
%                            move the \newcommand block to your preamble.
%
%  Companion notes:  docs/HAR-RV-methodology.md   (the HAR specification)
%                    docs/har_rv_methodology.tex
%  Verification:     docs/verify_target_alignment.py   (3/3 checks pass)
%  Code compared:    HAR-RV_RUN.PY
%                    exp/exp_long_term_forecasting.py
%                    data_provider/data_loader.py, utils/metrics.py, run.py
% =============================================================================
\documentclass[11pt,a4paper]{article}

\usepackage[T1]{fontenc}
\usepackage[utf8]{inputenc}
\usepackage{amsmath,amssymb,amsthm}
\usepackage{booktabs}
\usepackage{array}
\usepackage{bm}
\usepackage[margin=2.5cm]{geometry}
\usepackage[hidelinks]{hyperref}

\newcommand{\RV}{\mathrm{RV}}
\newcommand{\E}{\mathbb{E}}
\newcommand{\Var}{\operatorname{Var}}
\newcommand{\Fset}{\mathcal{F}}
\newcommand{\MSE}{\mathrm{MSE}}
\newcommand{\MAE}{\mathrm{MAE}}
\newcommand{\QLIKE}{\mathrm{QLIKE}}
\newcommand{\Lpath}{\mathcal{L}_{\mathrm{path}}}
\newcommand{\Lagg}{\mathcal{L}_{\mathrm{agg}}}
\newcommand{\Lin}{L_{\mathrm{in}}}
\newcommand{\Llab}{L_{\mathrm{lab}}}
\newcommand{\Lout}{L_{\mathrm{out}}}

\newtheorem{proposition}{Proposition}
\newtheorem{remark}{Remark}

\title{Forecast Targets and Loss Functions:\\
Deep Time-Series Models versus the Multi-Horizon HAR-RV Arrangement}
\date{}

\begin{document}
\maketitle

\begin{abstract}
\noindent
A deep forecasting model in the Time-Series-Library convention and the HAR-RV model of
Corsi (2009) do not, by default, predict the same object or minimise the same criterion.
The former predicts a \emph{path} of $\Lout$ daily values and minimises the mean squared
error averaged over that path; the latter predicts a \emph{scalar}, the $h$-day forward
average of realized variance, and minimises the squared error of that single quantity.
This note states both objects formally, proves the exact relationship between the two
losses --- path loss equals the HAR criterion plus a within-window dispersion term
(Proposition~\ref{prop:decomp}) --- documents the modifications that place the deep models
on the HAR target, and lists what remains non-comparable. All quantitative claims are
verified numerically in \texttt{docs/verify\_target\_alignment.py}.
\end{abstract}

%%% BODY BEGIN %%%

\section{Notation and the common sample}
\label{sec:setup}

Let $t=1,\dots,T$ index trading days, $\RV_t>0$ be daily realized variance and $\Fset_t$ the
information set at the close of day $t$. Both model families are estimated on the same
series after the same filtering rule (non-positive $\RV$ rows, i.e.\ non-trading days, are
dropped), and the chronological splits are
\begin{equation}
  \underbrace{\text{train} \le 2022}_{\text{deep: gradient updates}}
  \;\;\Big|\;\;
  \underbrace{2023\text{--}2024}_{\text{deep: early stopping}}
  \;\;\Big|\;\;
  \underbrace{\text{test} \ge 2025}_{\text{both: scored}},
  \qquad
  \underbrace{\text{train} \le 2024}_{\text{HAR: OLS}}
  \;\;\Big|\;\;
  \underbrace{\text{test} \ge 2025}_{\text{both: scored}} .
  \label{eq:splits}
\end{equation}
The HAR has no hyperparameters, so its estimation sample absorbs the validation years; the
\emph{test} window is identical for both, which is what makes the comparison legitimate.
Section~\ref{sec:caveats} returns to the fact that the two estimation samples are
nonetheless not the same.

Deep models are windowed by the triple $(\Lin,\Llab,\Lout)$: look-back (\texttt{seq\_len}),
decoder start-token length (\texttt{label\_len}) and forecast length (\texttt{pred\_len}).
Throughout we set $\Lout=h$ and consider the univariate case
(\texttt{features='S'}, one channel), which is the configuration used for realized variance.

\section{The default deep-learning arrangement}
\label{sec:dl-default}

\subsection{Target: a path}

For a forecast origin $t$ the network maps the look-back window to the \emph{whole future
path},
\begin{equation}
  \hat{\bm{y}}_{t} \;=\; g_\theta\bigl(\RV_{t-\Lin+1},\dots,\RV_{t}\bigr)
  \;\in\; \mathbb{R}^{\Lout},
  \qquad
  \bm{y}_{t} \;=\; \bigl(\RV_{t+1},\dots,\RV_{t+\Lout}\bigr)' .
  \label{eq:dl-target}
\end{equation}
The dependent variable is therefore $\Lout$-dimensional: the model is asked \emph{when},
not only \emph{how much}. The decoder start token consists of the last $\Llab$ observations
of the input window, all dated $\le t$, so \eqref{eq:dl-target} involves no look-ahead
despite appearances.\footnote{In the loader, $\texttt{batch\_y[:, :}\Llab\texttt{]}$ indexes
rows $[t-\Llab+1,\,t]$, which lie inside the encoder window; only the trailing $\Lout$
positions of the decoder input are zero-filled.}

\subsection{Loss: mean squared error averaged over the path}

Training minimises \texttt{nn.MSELoss()} over a minibatch $\mathcal{B}$,
\begin{equation}
  \Lpath(\theta)
  \;=\;
  \frac{1}{|\mathcal{B}|\,\Lout}
  \sum_{t\in\mathcal{B}}\;\sum_{k=1}^{\Lout}
  \bigl(\RV_{t+k}-\hat{y}_{t,k}\bigr)^{2},
  \label{eq:dl-loss}
\end{equation}
i.e.\ the squared error is averaged \emph{across horizon steps as well as across windows}.
Every step $k=1,\dots,\Lout$ carries equal weight, so a one-day-ahead error and a
twenty-two-day-ahead error contribute equally. Early stopping monitors the same functional
on the validation split.

\subsection{Reported metrics}

The harness reports $\MAE$, $\MSE$, $\mathrm{RMSE}$, $\mathrm{MAPE}$ and $\mathrm{MSPE}$
computed over the flattened $(n\times\Lout)$ tensor of path forecasts. Two of these deserve
a warning in any volatility application: $\mathrm{MAPE}$ and $\mathrm{MSPE}$ divide by
$\RV_t$ and therefore explode on quiet days, and $\MAE$ elicits the conditional
\emph{median} and is not robust to noise in the volatility proxy
(Section~\ref{sec:consistency}). Neither should drive model selection.

\subsection{Scaling}

The stock Time-Series-Library loader standardises each channel with statistics fitted on
the training split, $z=(y-\hat\mu_{\text{train}})/\hat\sigma_{\text{train}}$, and trains on
$z$. In this repository \texttt{Dataset\_Custom} is instantiated with \texttt{scale=False}
and \texttt{data\_factory} never overrides it, so \textbf{no standardisation is applied}:
the network consumes $\RV$ (or $\ln \RV$ under \texttt{-{}-log}) directly and
\eqref{eq:dl-loss} is expressed in variance units. Two consequences follow.
\begin{enumerate}
  \item The deep training loss and the HAR training loss are in the same units and are
        numerically comparable --- the reason for the choice.
  \item Raw $\RV$ is strongly right-skewed and badly conditioned as a network input. Under
        \texttt{-{}-log} the logarithm plays the variance-stabilising role that $z$-scoring
        would otherwise play; in raw mode nothing does, so raw-mode deep runs should be
        expected to optimise less comfortably than log-mode runs. This is a property of the
        experiment, not a defect, but it must be reported.
\end{enumerate}

\section{The HAR-RV target for comparison}
\label{sec:har}

The HAR predicts the scalar $h$-day forward \emph{average} realized variance
\begin{equation}
  Y^{(h)}_{t} \;=\; \frac{1}{h}\sum_{k=1}^{h}\RV_{t+k}
  \qquad\text{(raw)},
  \qquad
  y^{(h)}_{t} \;=\; \ln Y^{(h)}_{t}
  \qquad\text{(log-of-mean)},
  \label{eq:har-target}
\end{equation}
by ordinary least squares on the three heterogeneous components, so its in-sample criterion
is the squared error of a \emph{single} number per origin. Full specification: see the
companion note \texttt{docs/HAR-RV-methodology.md}. The essential point here is that
\eqref{eq:dl-target} and \eqref{eq:har-target} are different objects whenever $h>1$:
one is a vector in $\mathbb{R}^{h}$, the other its arithmetic mean.

\section{The modifications: \texttt{-{}-aggregate\_mean} and \texttt{-{}-log}}
\label{sec:modifications}

Six changes place the deep models on the HAR target. Each is stated with the mathematics it
implements.

\paragraph{(i) Scalar output head.} The model is constructed with
$\texttt{pred\_len}=1$, so the head emits $\hat{Y}^{(h)}_{t}\in\mathbb{R}$ rather than a
path, while the loader still returns the full $h$-day future window needed to build the
ground truth. The forecast is a direct projection onto the aggregate, not a post-hoc
average of a path forecast.

\paragraph{(ii) Target aggregation in variance space.} The future window is reduced to
\eqref{eq:har-target}. Aggregation always happens on the variance scale --- the additive
scale, where averaging variances is the meaningful operation --- and the logarithm, when
present, is applied to the aggregate:
\begin{equation}
  Y^{(h)}_{t}=\frac{1}{h}\sum_{k=1}^{h}\RV_{t+k},
  \qquad
  y^{(h)}_{t}=\ln\Bigl(\frac{1}{h}\sum_{k=1}^{h}\RV_{t+k}\Bigr)
  \;\ne\;
  \frac{1}{h}\sum_{k=1}^{h}\ln \RV_{t+k}.
  \label{eq:logofmean}
\end{equation}
The right-hand inequality is Jensen's: the mean of logs is the log of the \emph{geometric}
mean, a smaller and mechanically smoother object. Only log-of-mean gives
$\exp\bigl(y^{(h)}_{t}\bigr)=Y^{(h)}_{t}$ exactly at every horizon, which is what allows a
log-scale run to be scored against a raw-scale run and against the HAR.

\paragraph{(iii) Numerically stable log target.} Under \texttt{-{}-log} the stored series is
already $x_s=\ln \RV_s$, so the target is formed without ever exponentiating a large
magnitude:
\begin{equation}
  y^{(h)}_{t}
  \;=\;\ln\Bigl(\frac{1}{h}\sum_{k=1}^{h}e^{x_{t+k}}\Bigr)
  \;=\;\operatorname*{logsumexp}_{k=1,\dots,h}\bigl(x_{t+k}\bigr)-\ln h ,
  \label{eq:logsumexp}
\end{equation}
an algebraic identity, verified to $8.9\times10^{-16}$ on the sample.

\paragraph{(iv) Aggregate loss.} The training criterion becomes
\begin{equation}
  \Lagg(\theta)
  \;=\;
  \frac{1}{|\mathcal{B}|}\sum_{t\in\mathcal{B}}
  \bigl(Y^{(h)}_{t}-\hat{Y}^{(h)}_{t}\bigr)^{2}
  \qquad\text{(raw)},
  \qquad
  \frac{1}{|\mathcal{B}|}\sum_{t\in\mathcal{B}}
  \bigl(y^{(h)}_{t}-\hat{y}^{(h)}_{t}\bigr)^{2}
  \qquad\text{(log)},
  \label{eq:agg-loss}
\end{equation}
which is precisely the criterion OLS minimises in the HAR regression. The two model
families now optimise the same functional of the same target, differing only in the
hypothesis class.

\paragraph{(v) No inverse transform in aggregate mode.} With standardisation off there is
nothing to invert; and even were it on, the aggregated target does not live in the
scaler's space, since $\mathcal{S}^{-1}$ applied to an average of standardised variances is
not the average of the inverse-transformed variances unless $\mathcal{S}$ is affine and
applied before aggregation. Aggregate-mode outputs are therefore left on the modelling
scale and scored there.

\paragraph{(vi) Evaluation on the HAR footing.} Beyond the harness metrics, aggregate-mode
runs report the HAR loss set: $\MSE$, $\MAE$ and $\QLIKE$ (Patton, 2011),
\begin{equation}
  \QLIKE=\frac{1}{n}\sum_{t}\left[\frac{Y^{(h)}_{t}}{\hat{Y}^{(h)}_{t}}
  -\ln\!\left(\frac{Y^{(h)}_{t}}{\hat{Y}^{(h)}_{t}}\right)-1\right]\;\ge\;0,
  \label{eq:qlike}
\end{equation}
with non-positive forecasts floored at $10^{-4}\,\overline{\RV}_{\text{train}}$ and the
count reported rather than discarded --- the identical rule to the HAR script, so the floor
cannot silently advantage either model. Under \texttt{-{}-log} both sides are mapped back to
variances first (Section~\ref{sec:backtransform}).

\paragraph{Sample hygiene.} \texttt{-{}-aggregate\_mean} and \texttt{-{}-log} both force
\texttt{drop\_nonpositive}, so the deep models train on exactly the rows the HAR keeps; and
the test loader sets \texttt{drop\_last=False}, so the trailing partial batch is retained.
Dropping it would discard up to $|\mathcal{B}|-1$ test targets and break the row-for-row
correspondence established in Proposition~\ref{prop:align}.

\section{The exact relationship between the two losses}
\label{sec:decomposition}

\begin{proposition}[Path loss $=$ aggregate loss $+$ within-window dispersion]
\label{prop:decomp}
Fix an origin $t$ and a horizon $h$. Let $e_k=\RV_{t+k}-\hat{y}_{t,k}$ be the daily forecast
errors of a path forecast, $k=1,\dots,h$, and let
$\bar{e}=h^{-1}\sum_{k=1}^{h}e_k$. Then
\begin{equation}
  \underbrace{\frac{1}{h}\sum_{k=1}^{h}e_k^{2}}_{\text{path loss at }t}
  \;=\;
  \underbrace{\bar{e}^{\,2}}_{\text{aggregate (HAR) loss at }t}
  \;+\;
  \underbrace{\frac{1}{h}\sum_{k=1}^{h}\bigl(e_k-\bar{e}\bigr)^{2}}_{\text{within-window dispersion}\;\ge\,0},
  \qquad
  \bar{e}=Y^{(h)}_{t}-\frac{1}{h}\sum_{k=1}^{h}\hat{y}_{t,k}.
  \label{eq:decomp}
\end{equation}
\end{proposition}

\begin{proof}
The first equality is the algebraic identity
$h^{-1}\sum_k e_k^2 = \bar{e}^{\,2}+h^{-1}\sum_k(e_k-\bar{e})^2$, valid for any finite
sequence. The expression for $\bar{e}$ follows from linearity of the average and the
definition \eqref{eq:har-target} of $Y^{(h)}_{t}$.
\end{proof}

Three consequences, each of which should be stated when the two model families are compared.

\begin{enumerate}
  \item \textbf{The default deep loss is not the HAR criterion; it is the HAR criterion plus
        a penalty.} Since the dispersion term is non-negative,
        $\Lpath\ge\Lagg$ always, with equality only when the error is constant across the
        window. A model trained under \eqref{eq:dl-loss} spends capacity on getting the
        \emph{timing} of volatility within the window right --- something the HAR is never
        asked to do and is never scored on.
  \item \textbf{The penalty dominates in practice.} On the shipped EUR/USD sample, using a
        plausible random-walk-with-noise path forecaster, the median ratio
        $\Lpath/\Lagg$ is $2.8$ at $h=5$ and $4.4$ at $h=22$, so the dispersion term
        accounts for a median $65\%$ of the path loss at $h=5$ and $77\%$ at $h=22$. Two
        thirds to three quarters of the default deep training signal is within-window
        \emph{shape}, not the level of the $h$-day average, and the share grows with the
        horizon. Comparing such a model's loss with the HAR's is comparing two different
        quantities.
  \item \textbf{Aggregating after the fact is not the same as training on the aggregate.}
        Averaging a trained path forecast gives a valid $\hat{Y}^{(h)}$, but the parameters
        were selected to minimise $\Lagg + \text{dispersion}$, not $\Lagg$. Modification (i)
        removes the second term from the objective rather than from the report.
\end{enumerate}

\section{Which loss elicits which functional}
\label{sec:consistency}

A loss $L$ is \emph{consistent} for a functional $\Gamma$ if
$\Gamma(\Fset_t)=\arg\min_{f}\E[L(Y,f)\mid\Fset_t]$. This determines what a trained model
converges to and is the criterion by which a training loss and an evaluation loss can be
said to agree.

\begin{center}
\begin{tabular}{@{}llp{5.6cm}@{}}
\toprule
Loss & Elicits & Comment \\
\midrule
Path $\MSE$ \eqref{eq:dl-loss}      & $\E[\RV_{t+k}\mid\Fset_t]$, each $k$ & the whole conditional mean path \\
Aggregate $\MSE$ \eqref{eq:agg-loss}, raw & $\E[Y^{(h)}_{t}\mid\Fset_t]$   & the HAR estimand \\
Aggregate $\MSE$, log scale         & $\E[\ln Y^{(h)}_{t}\mid\Fset_t]$     & \emph{not} $\E[Y^{(h)}_t\mid\Fset_t]$; needs \S\ref{sec:backtransform} \\
$\QLIKE$ \eqref{eq:qlike}           & $\E[Y^{(h)}_{t}\mid\Fset_t]$         & same functional as raw $\MSE$ \\
$\MAE$                              & conditional median                   & different functional; not proxy-robust \\
\bottomrule
\end{tabular}
\end{center}

Two readings follow. First, \emph{training on $\MSE$ and evaluating on $\QLIKE$ is
coherent}: both belong to Patton's (2011) robust family, both are minimised by the
conditional mean of the variance proxy, and they differ only in how they weight errors
across the volatility range --- $\QLIKE$ penalising under-prediction far more heavily.
Neither model optimises $\QLIKE$ directly, which is worth stating, but the ranking it
produces is not out-of-objective in the way an $\MAE$ ranking would be. Second, $\MAE$ and
$\mathrm{MAPE}$ elicit different functionals and are not robust to the noise in $\RV$ as a
proxy for latent variance; report them descriptively, select on $\MSE$ and $\QLIKE$.

\section{Back-transformation under \texttt{-{}-log}}
\label{sec:backtransform}

Log-scale training elicits $\E[\ln Y^{(h)}\mid\Fset_t]$, so $\exp(\hat{y}^{(h)})$ is a
conditional \emph{median} and understates the conditional mean. Both model families apply
the lognormal correction, but the deep version carries one extra term:
\begin{equation}
  \hat{Y}^{(h)}_{t}
  \;=\;
  \exp\!\Bigl(\hat{y}^{(h)}_{t} \;+\; \underbrace{b}_{\text{mean residual}} \;+\; \tfrac{1}{2}\hat\sigma_u^{2}\Bigr),
  \qquad
  b=\frac{1}{n_{\text{tr}}}\sum_{t\in\text{train}}\hat{u}_t,
  \qquad
  \hat\sigma_u^{2}=\widehat{\Var}_{\text{train}}(\hat{u}_t).
  \label{eq:backtransform}
\end{equation}
For OLS with an intercept the mean residual is \emph{exactly} zero, so $b=0$ and
\eqref{eq:backtransform} reduces to the HAR correction $\exp(\hat y+\hat\sigma_u^2/2)$; a
network has no such guarantee, and omitting $b$ would leave a systematic level error in the
back-transformed forecast. Setting $b=\hat\sigma_u^2=0$ recovers the naive transform, which
is reported alongside so the size of the correction stays visible.

Both moments are computed on \textbf{training} residuals only. Using test residuals would
leak the realised outcome into the forecast. The result is strictly positive, so a
log-scale model cannot produce the non-positive variance forecasts an unconstrained
raw-scale head can.

\section{Test-sample alignment}
\label{sec:alignment}

\begin{proposition}[The two families score the same forecast origins]
\label{prop:align}
Let $N$ be the number of retained rows and $V=\#\{s:\mathrm{year}(s)\le 2024\}$. With
$\Lout=h$, the deep loader yields
\begin{equation}
  n_{\text{DL}} \;=\; \bigl(N-(V-\Lin)\bigr)-\Lin-h+1 \;=\; N-V-h+1
  \label{eq:dl-count}
\end{equation}
test windows, whose targets span days $[V+i,\,V+i+h)$ for $i=0,\dots,n_{\text{DL}}-1$;
the HAR retains rows $t\in[V,\,N-h]$, i.e.\ $n_{\text{HAR}}=N-V-h+1$ targets spanning
$[t,\,t+h)$. The two sets of forecast origins, and hence the two sets of scored dates,
coincide exactly, independently of $\Lin$.
\end{proposition}

The cancellation of $\Lin$ in \eqref{eq:dl-count} is not accidental: the validation and test
windows begin $\Lin$ rows early precisely so that the first origin in each split has a
complete look-back window without consuming a target. Numerically, on the shipped sample
($T=3766$, $\Lin=96$):

\begin{center}
\begin{tabular}{@{}llll@{}}
\toprule
$h$ & $n_{\text{DL}}$ & $n_{\text{HAR}}$ & scored window \\
\midrule
$1$  & $402$ & $402$ & 2025-01-02 -- 2026-07-21 \\
$5$  & $398$ & $398$ & 2025-01-02 -- 2026-07-15 \\
$22$ & $381$ & $381$ & 2025-01-02 -- 2026-06-22 \\
\bottomrule
\end{tabular}
\end{center}

Losses are therefore averages over identical index sets and can be placed in the same table
without reweighting. Note that the test targets overlap for $h>1$, so any formal comparison
--- a Diebold--Mariano test on the loss differential, say --- must use a HAC variance with
bandwidth at least $h-1$; the conventional choice, and the one the HAR script uses, is
$\max\{2(h-1),\lfloor 4(n/100)^{2/9}\rfloor\}$.

\section{Side-by-side summary}
\label{sec:summary}

\begin{center}
\small
\begin{tabular}{@{}p{3.1cm}p{3.9cm}p{3.9cm}p{3.4cm}@{}}
\toprule
 & \textbf{Deep, default} & \textbf{Deep, \texttt{-{}-aggregate\_mean}} & \textbf{HAR-RV} \\
\midrule
Target
 & path $(\RV_{t+1},\dots,\RV_{t+h})'$
 & scalar $Y^{(h)}_t=h^{-1}\sum_k \RV_{t+k}$
 & scalar $Y^{(h)}_t$ \\[2pt]
Output dimension
 & $h$
 & $1$
 & $1$ \\[2pt]
Training loss
 & $\Lpath$, eq.\ \eqref{eq:dl-loss}
 & $\Lagg$, eq.\ \eqref{eq:agg-loss}
 & $\Lagg$ (OLS) \\[2pt]
Elicits
 & conditional mean path
 & $\E[Y^{(h)}_t\mid\Fset_t]$
 & $\E[Y^{(h)}_t\mid\Fset_t]$ \\[2pt]
Conditioning set
 & $\Lin$ lags ($96$)
 & $\Lin$ lags ($96$)
 & $22$ lags, 3 aggregates \\[2pt]
Parameters
 & $10^{4}$--$10^{6}$
 & $10^{4}$--$10^{6}$
 & $4$ \\[2pt]
Log mode
 & $\ln \RV$ per day
 & $\ln$ of the mean, eq.\ \eqref{eq:logsumexp}
 & $\ln$ of the mean \\[2pt]
Standardisation
 & off (\texttt{scale=False})
 & off
 & none \\[2pt]
Back-transform
 & ---
 & $\exp(\hat y+b+\hat\sigma_u^2/2)$
 & $\exp(\hat y+\hat\sigma_u^2/2)$, $b\equiv0$ \\[2pt]
Reported metrics
 & MAE, MSE, RMSE, MAPE, MSPE
 & $+$ QLIKE, MSE\_RV, MAE\_RV, naive QLIKE
 & MSE, MAE, QLIKE $(+$ RV-scale$)$ \\[2pt]
Estimation sample
 & $\le 2022$ $(+$ val for stopping$)$
 & $\le 2022$ $(+$ val for stopping$)$
 & $\le 2024$ \\[2pt]
Inference
 & none
 & none
 & Newey--West HAC \\
\bottomrule
\end{tabular}
\end{center}

\section{What remains non-comparable}
\label{sec:caveats}

Aligning the target and the loss removes the largest source of incomparability but not all
of it. The following differences survive and belong in the limitations paragraph.

\begin{enumerate}
  \item \textbf{Conditioning sets differ.} The HAR conditions on $22$ lags compressed into
        three averages; the deep models condition on $\Lin=96$ raw lags. The deep models
        therefore see more history \emph{and} a richer representation of it. For a strict
        like-for-like test, run the deep models with $\Lin=22$; otherwise report the
        asymmetry, since part of any deep-model advantage is simply a longer window.
  \item \textbf{Estimation samples differ.} By \eqref{eq:splits} the HAR fits on data
        through 2024 while the deep models take gradient steps only on data through 2022,
        using 2023--2024 for early stopping. The HAR has roughly two additional years of
        fitting data; the deep models use those years for model selection instead. This is
        the standard arrangement --- OLS has nothing to tune --- but it is not an identical
        information budget.
  \item \textbf{Neither family optimises $\QLIKE$.} Both train on squared error and are
        ranked partly on $\QLIKE$. Section~\ref{sec:consistency} argues this is coherent,
        since both losses elicit the same functional, but a $\QLIKE$-trained variant is the
        natural robustness check.
  \item \textbf{Model selection is stochastic on one side only.} Deep results depend on the
        seed, initialisation and early-stopping epoch; OLS is deterministic and has a closed
        form. Report deep results over several seeds, with dispersion, not a single run.
  \item \textbf{Uncertainty quantification.} The HAR delivers HAC standard errors and
        testable coefficient restrictions; the deep models as configured deliver point
        forecasts only.
  \item \textbf{The Jensen correction is estimated per model.} Under \texttt{-{}-log} each
        model carries its own $(b,\hat\sigma_u^2)$, so back-transformed comparisons include
        a model-specific constant. Report the correction factor $\exp(b+\hat\sigma_u^2/2)$
        alongside the losses, as the harness does.
\end{enumerate}

\section{Reporting recipe}
\label{sec:recipe}

For each horizon $h\in\{1,5,22\}$ and each scale (raw, log), run the deep models with
$\Lout=h$ and \texttt{-{}-aggregate\_mean}, and the HAR at the same $h$; then report, on the
identical test index set of Proposition~\ref{prop:align}:
\begin{enumerate}
  \item $\MSE$ and $\MAE$ on the modelling scale, with the scale named in the column header
        (raw $\RV$ versus $\ln \RV$ units --- these are not interchangeable);
  \item $\QLIKE$ on the variance scale, with the naive-transform value and the count of
        floored forecasts in a footnote;
  \item $\MSE_{\RV}$ and $\MAE_{\RV}$ after back-transformation, which are the only
        cross-scale-comparable columns;
  \item Diebold--Mariano statistics against the HAR benchmark with a HAC variance at
        bandwidth $\ge h-1$, or a Model Confidence Set if more than two models compete;
  \item for deep models, mean and standard deviation across seeds.
\end{enumerate}
State in the caption that the deep target is $Y^{(h)}_t$ rather than the daily path, and
cite Proposition~\ref{prop:decomp} for why the distinction matters --- a reader who assumes
the default path loss will otherwise misread every deep column in the table.

%%% BODY END %%%

\begin{thebibliography}{9}
\bibitem{corsi2009} Corsi, F. (2009). A simple approximate long-memory model of realized volatility. \emph{Journal of Financial Econometrics}, 7(2), 174--196.
\bibitem{dm1995} Diebold, F.~X., \& Mariano, R.~S. (1995). Comparing predictive accuracy. \emph{Journal of Business \& Economic Statistics}, 13(3), 253--263.
\bibitem{gneiting2011} Gneiting, T. (2011). Making and evaluating point forecasts. \emph{Journal of the American Statistical Association}, 106(494), 746--762.
\bibitem{hln2011} Hansen, P.~R., Lunde, A., \& Nason, J.~M. (2011). The model confidence set. \emph{Econometrica}, 79(2), 453--497.
\bibitem{patton2011} Patton, A.~J. (2011). Volatility forecast comparison using imperfect volatility proxies. \emph{Journal of Econometrics}, 160(1), 246--256.
\bibitem{ps2009} Patton, A.~J., \& Sheppard, K. (2009). Evaluating volatility and correlation forecasts. In T.~G. Andersen et al. (Eds.), \emph{Handbook of Financial Time Series} (pp.~801--838). Springer.
\bibitem{abdl2003} Andersen, T.~G., Bollerslev, T., Diebold, F.~X., \& Labys, P. (2003). Modeling and forecasting realized volatility. \emph{Econometrica}, 71(2), 579--625.
\bibitem{duan1983} Duan, N. (1983). Smearing estimate: A nonparametric retransformation method. \emph{Journal of the American Statistical Association}, 78(383), 605--610.
\end{thebibliography}

\end{document}
